\section{Gaussian scales}\label{sec:gaussian}

We next identify the fluctuation law on scales where the independent
common-word increments have negligible large jumps. There are three
steps: an annealed limit for independent paths, a comparison for shared
paths away from one another, and an occupation estimate excluding a
singular interaction on their projected diagonal. The resulting shared
two-copy limit will imply a quenched limit in probability.

\begin{definition}[Gaussian sequence]\label{gauss-sequence}
A sequence \(r\to\infty\) is \emph{Gaussian} if
\[
n(r)\Pr(|S|>a r)\longrightarrow0\qquad\text{for every fixed }a>0.
\tag{14}\label{eq:14}
\]
\end{definition}

Throughout this section limits are along any fixed Gaussian sequence.
Height-indexed processes are linearly interpolated on the grid
\(h/n(r)\). Functional convergence means weak convergence in
\(C([0,T],\mathbb R^k)\), with the uniform norm, for every fixed finite
\(T\); a slightly larger horizon is used when interpolating near its
endpoint. Put \(\sigma_*^2=1/(2\bar L)\).

\subsection{Independent-path Brownian limit}

The common-word increments give a Brownian limit at common cuts. To
obtain a limit for a single walk at every first-hit level, we must also
control the fluctuations between cuts and recover the deterministic
centering. Neither step requires a second moment for word decorations.

\begin{lemma}[Independent first-hit fluctuations]\label{gauss-independent}
For two independent conditioned paths, \(\Delta_h/r\) converges
functionally to Brownian motion of variance parameter \(1/\bar L\).
For a single \(P_0^+\) path, \((Y_h-b_h)/r\) converges functionally to
Brownian motion of variance parameter \(\sigma_*^2\).
\end{lemma}

\begin{proof}
Write \((L_i,S_i)\) for the successive independent common-word widths
and difference increments, and put
\[
 C_m=\sum_{i=1}^mL_i,\qquad A_m=\sum_{i=1}^mS_i,
 \qquad N(h)=\max\{m:C_m\le h\}.
\]
The difference at the preceding common cut is
\(\widehat\Delta_h=A_{N(h)}\). As with the first-hit process, we
interpolate its values linearly on the height grid \(h/n(r)\).

\emph{The limit at common cuts.}
First consider \(A_m/r\), with \(n(r)\) summands per unit time.
The Gaussian-sequence assumption permits a symmetric truncation
\[
 S_i^{(r)}=S_i\mathbf1_{\{|S_i|\le\eta_r r\}},
 \qquad \eta_r\downarrow0,
\]
for which
\[
 n(r)\Pr(|S|>\eta_r r)\longrightarrow0,
 \qquad
 \frac{\mathbb E[(S_i^{(r)})^2]}{m(r)}\longrightarrow1.
\]
Indeed, for every fixed \(0<\eta<1\),
\[
 0\le m(r)-\mathbb E[S^2\mathbf1_{\{|S|\le\eta r\}}]
 \le r^2\Pr(|S|>\eta r)=o(m(r));
\]
a sufficiently slow diagonal choice of \(\eta_r\) gives both
requirements. The probability that any summand is removed on a fixed
time interval tends to zero.

The variables \(S_i^{(r)}/r\) have mean zero, variance
\((1+o(1))/n(r)\), and fourth moment at most \(\eta_r^2\) times
that variance. Taylor expansion of their characteristic functions
therefore gives the standard Brownian finite-dimensional laws. For
tightness, the fourth-moment computation for independent centered sums,
followed by the fourth-moment maximal inequality, bounds the expected
maximal fourth power on a block of \(\lfloor\delta n(r)\rfloor\)
summands by
\[
 O(\delta^2+\delta\eta_r^2).
\]
A union bound over the blocks in a fixed time interval, followed by
\(\delta\downarrow0\), proves tightness. Thus \(A_m/r\) has the
standard Brownian functional limit.

The width strong law gives
\[
 \sup_{0\le h\le Tn(r)}
 \left|\frac{N(h)}{n(r)}-\frac{h}{\bar L n(r)}\right|
 \longrightarrow0
 \quad\text{in probability}.
\]
Applying this time change to the tight sum process, whose limits are
continuous, proves the Brownian limit for
\(\widehat\Delta_h/r\), with variance parameter \(1/\bar L\).
The counts and the sums need not be independent. The maximal normalized
jump tends to zero as well, so the stated interpolation convention has
the same limit.

\emph{Passing from cuts to every first hit.}
Lemma~\ref{fluct-local-tightness} makes
\(\Delta_h-\widehat\Delta_h\) tight uniformly at deterministic
levels \(h\). This already gives the required finite-dimensional
limits. To prove path tightness, fix \(\delta>0\) and partition the
height range into consecutive blocks of
\(\lfloor\delta n(r)\rfloor\) levels. There are only finitely many
block starts for fixed \(T,\delta\).

Suppose that within one of these blocks, with start \(a\),
\[
 |\Delta_h-\Delta_a|>\rho r
\]
for some level \(h\) in the block. Under the raw independent pair
law, select the first such offending level among prefixes for which
both walks have avoided the initial drop. The selection is determined
by the paired prefixes through that level. The raw probability of
finding an offense is at least \(p^2\) times its probability under
the independent conditioned pair law. Imposing both fresh truths at
the selected level multiplies each selected prefix weight by \(p^2\),
leaves the offense unchanged, and implies both initial no-drop events.
After division by the initial normalizer \(p^2\), both probabilities
below are under the conditioned pair law, and
\[
 \Pr(\text{an offense})
 \le p^{-2}\Pr(\text{an offense at a common cut}).
\]
The same statement holds for the union of the blocks by selecting its
first offense.

At a common cut \(h\), the offending difference satisfies
\[
 |\Delta_h-\Delta_a|
 \le |\widehat\Delta_h-\widehat\Delta_a|
       +|\Delta_a-\widehat\Delta_a|.
\]
The maximum of the second term over the finitely many deterministic
block starts is \(o_{\Pr}(r)\). The first term is controlled by the
modulus of the common-cut process, which has a continuous Brownian
limit. Consequently, for each \(\rho>0\), the limsup probability of
an offense tends to zero as \(\delta\downarrow0\). Adjacent-block
bounds and interpolation give tightness of \(\Delta_h/r\), proving
the first assertion.

\emph{Recovering a single path.}
Let \(\mu_r\) be the law in continuous path space of the interpolated
single-path process \(Y_h/r\), on a fixed compact time interval. Its
iid differences are tight by the first assertion. We first show that
\(\mu_r\) is tight after suitable deterministic path translations.
Choose compact sets \(K_j\) capturing the difference laws with errors
\(\varepsilon_j^2\), where \(\varepsilon_j\downarrow0\) and
\(\sum_j\varepsilon_j<1\). If a prospective center \(f\) is sampled
from \(\mu_r\), then
\[
 \int\mu_r(f+K_j)\,\mu_r(df)\ge1-\varepsilon_j^2.
\]
Markov's inequality shows that the centers for which
\(\mu_r(f+K_j)<1-\varepsilon_j\) have \(\mu_r\)-probability at
most \(\varepsilon_j\). Some center path \(f_r\) therefore works
for every \(j\), and the translated laws are tight. We may choose
\(f_r\) with the same interpolation as the original paths.

Along a subsequence on which these translated laws converge, let
\(G\) be the continuous limiting process. Its iid two-copy difference
has the Brownian law just proved. Cram\'er's normal decomposition
principle~\cite{Cramer1936}, in the iid-difference form verified below,
applied to every scalar linear combination of finitely many coordinates
of \(G\), shows that \(G\) has Gaussian finite-dimensional laws.
Its covariance is half that of the difference process, namely
\[
 \operatorname{Cov}(G(s),G(t))=\sigma_*^2\min(s,t).
\]
Only a deterministic mean function remains to be identified.

For completeness, suppose that \(Z-Z'\) is normal for iid real
variables \(Z,Z'\). Choose \(C_0\) with
\(\Pr(|Z'|\le C_0)>0\). Intersecting either tail of \(Z\) with
this event bounds that tail by a constant times a Gaussian tail. Hence
\(f(z)=\mathbb E e^{zZ}\) is entire and obeys
\(|f(z)|\le\exp(C(1+|z|^2))\). The identity
\(f(z)f(-z)=\mathbb E e^{z(Z-Z')}\), first on the real axis and
then everywhere, shows that \(f\) has no zeros. Its entire logarithm
\(g\), chosen with \(g(0)=0\), satisfies
\(\operatorname{Re}g(z)\le C(1+|z|^2)\). On every circle the mean
of \(\operatorname{Re}g\) is zero, so its absolute integral mean is
also \(O(1+|z|^2)\). The Fourier coefficient formula on circles of
arbitrarily large radius then makes every Taylor coefficient of \(g\)
of order greater than two vanish. Thus \(g\) is quadratic, real on
the real axis, with \(g''(0)=\operatorname{Var}Z\ge0\). It follows
that \(Z\) is normal, allowing a point mass.

Write \(m_G(t)=\mathbb E G(t)\). Sample continuity and convergence
in law of Gaussian marginals at convergent times imply that \(m_G\)
is continuous. Each marginal has the unique median \(m_G(t)\),
including at time zero. At a grid time \(h/n(r)\), the number
\[
 \frac{b_h}{r}-f_r(h/n(r))
\]
is a median of the translated marginal. Along every sequence of grid
times converging to \(t\), weak path convergence and uniqueness of
the limiting median make these medians converge to \(m_G(t)\).
Compactness of the time interval makes the convergence uniform over
grid times; interpolation gives uniform convergence on the whole
interval. Use a slightly longer interval when interpolating at the
endpoint. Subtracting the interpolated medians from the original
process therefore removes precisely the limiting mean function.
Every subsequential limit is Brownian motion of variance parameter
\(\sigma_*^2\), proving the single-path assertion.
\end{proof}

\begin{lemma}[Linearization of the medians]\label{gauss-medians}
Uniformly for \(k+j\le T n(r)\), for each fixed \(T\),
\[
 b_{k+j}-b_k-b_j=o(r).
 \tag{15}\label{eq:15}
\]
Writing \(N_r=\lfloor n(r)\rfloor\) and
\(\theta_r=b_{N_r}/N_r\), we have
\[
 \max_{h\le Tn(r)}|b_h-h\theta_r|=o(r),
 \qquad \theta_r=o(r).
\]
\end{lemma}

\begin{proof}
We first prove the additive relation. It suffices to consider arbitrary
sequences of indices for which
\(k/n(r)\to s\) and \(j/n(r)\to t\), and to pass further so
that \(k\) is fixed or tends to infinity. Let \(\mathcal C_k\)
be exact single truth at level \(k\). Under \(P_0^+\), this event
has probability \(u_k\ge p\). Conditional on \(\mathcal C_k\),
the prefix has the raw law conditioned on \(T_k<T_{-1}\), and its
relative suffix is an independent fresh \(P_0^+\) path.

If \(k\to\infty\), the total variation distance between this
prefix law and the \(P_0^+\) prefix law is at most
\(1-p/u_k\to0\). If \(k\) is fixed, its displacement and median
are fixed random and deterministic quantities, respectively, and their
normalized difference tends to zero. Lemma~\ref{gauss-independent}
therefore gives, under the conditional law,
\[
 \frac{Y_{k+j}-b_k-b_j}{r}
 \ \Longrightarrow\ N(0,\sigma_*^2(s+t)).
\]
The same lemma gives this normal limit under the unconditional
\(P_0^+\) law when the centering is \(b_{k+j}\).

Set \(d_r=(b_k+b_j-b_{k+j})/r\). The unconditional distribution
centered by \(b_{k+j}\) dominates \(p\) times the conditional
distribution with that same centering. Tightness first forces \(d_r\)
to be bounded. Along a further subsequence with \(d_r\to d\), the
limiting domination would read
\[
 N(0,v)\ \ge\ pN(d,v),\qquad v=\sigma_*^2(s+t),
\]
as measures. For \(v>0\), a nonzero shift makes the ratio of the
shifted normal density to the unshifted one unbounded. For \(v=0\),
the two point masses have different supports unless \(d=0\).
Thus \(d=0\) in every case. Compactness of the ranges of
\(k/n(r)\) and \(j/n(r)\) proves the uniform relation
\eqref{eq:15}.

To turn approximate additivity into a linear approximation, put
\(N=N_r\), \(\theta=b_N/N\), and
\(d_j=b_j-j\theta\) for \(0\le j<N\). Apply \eqref{eq:15}
to \(j+j\), and also to \(N+(2j-N)\) when \(2j\ge N\).
Uniformly over these indices,
\[
 2d_j=d_{2j\bmod N}+o(r).
\]
The maximum of \(|d_j|\) is therefore at most half itself plus
\(o(r)\), so it is \(o(r)\). For \(h\le Tn(r)\), write
\(h=qN+j\), where \(q\) remains bounded, and apply
\eqref{eq:15} a bounded number of times. This proves
\(\max_{h\le Tn(r)}|b_h-h\theta_r|=o(r)\). Finally, at \(h=1\)
this gives \(|b_1-\theta_r|=o(r)\); since \(b_1\) is fixed,
\(\theta_r=o(r)\).
\end{proof}

\subsection{Shared paths: marginals and separation}

Two paths in one environment interact when their traces visit a common
site. We first show that the conditioning on the second path does not
change the limiting fluctuations of the first. We then bound trace
contact by the probabilities of missing two disjoint endpoint windows.
This comparison concerns the entire traces, including their excursions
between first hits of successive layers.

\begin{lemma}[Shared marginal limits]\label{gauss-shared-marginals}
Under \(P_{x,y}^{++}\), each marginal process
\[
B_i^r(h/n(r))=(Y_h^{(i)}-h\theta_r)/r
\tag{16}\label{eq:16}
\]
converges functionally to Brownian motion of variance parameter
\(\sigma_*^2\), uniformly over starts on a common layer. Thus the
conclusion holds along every choice \(x=x(r),y=y(r)\) with
\(x_1=y_1\), whatever their separation. The pair is jointly tight.
\end{lemma}

\begin{proof}
Relative to the single conditioned environment-and-path law from \(x\),
the first marginal has density
\[
\frac{p\,q(y)}{p_2(x,y)}.
\]
For a fixed integer \(J\), replace \(q(y)\) by \(a_J(y)\), then
normalize. The resulting probability law differs from the original
one in total variation by at most \(C(u_J-p)\), uniformly in \(x,y\).
Indeed \(a_J\ge q\), and the added mass before normalization is
\[
\frac{E[q(x)(a_J(y)-q(y))]}{p_2(x,y)}
\le\frac{u_J-p}{\underline p_2}.
\]
All these normalized laws have density bounded by a common constant
relative to the single conditioned law.

The weight \(a_J(y)\) uses only rows between the common base and the
layer \(J\) above it, excluding the latter. Stop the single-word chain
of the first walk at its first true boundary at or above that layer.
Every row used by the weight lies below this boundary. Exact-word
factorization therefore gives a fresh \(P^+\) suffix, independent of
the removed path prefix under the weighted law. The weight may depend
on infinitely many sites in those \(J\) layers; only their heights
matter for this factorization.

For fixed \(J\), the height and size of this initial prefix are tight
under \(P_x^+\), and hence under every one of the bounded weighted laws.
Its height divided by \(n(r)\) tends to zero. Its displacement divided
by \(r\) also tends to zero, as does its centering correction, since
\(\theta_r/r\to0\). The suffix limit and its tight continuous modulus
thus give the asserted limit under the approximating law. Send
\(r\to\infty\) first and then \(J\to\infty\). The same proof applies
to the second marginal, and tightness of the two marginals gives joint
tightness.
\end{proof}

The next elementary observation converts endpoint concentration into a
bound on trace contact. Fix an environment and a base layer of height
\(F\). Let each walk be conditioned never to fall below \(F\), and
stop it on first arrival at height \(F+k\). For disjoint sets
\(I_1,I_2\) in that top layer, write \(E_i\) for the event that walk
\(i\)'s endpoint lies outside \(I_i\). Then
\[
P_{x,\omega}(\,\cdot\mid D)\otimes
P_{y,\omega}(\,\cdot\mid D)
 \{\text{the two stopped traces meet}\}
\le 2P_{x,\omega}(E_1\mid D)+2P_{y,\omega}(E_2\mid D).
\]
Here the two occurrences of \(D\) use the same absolute floor \(F\).
To prove the bound, for \(F\le z_1\le F+k\) let \(v_i(z)\) be the probability of a top endpoint
in \(I_i\) for a quenched walk from \(z\) conditioned to avoid that
floor. We consider only sites with positive probability of such
avoidance. The Markov property makes \(v_i(z)\) the conditional endpoint
probability whenever either walk first visits \(z\), regardless of
which walk visits it. Since \(I_1,I_2\) are disjoint,
\(v_1(z)+v_2(z)\le1\).

Put \(A_i=\{z:v_i(z)\le1/2\}\). At the first visit to \(A_i\) before
or at the top, the conditional probability of \(E_i\) is at least
\(1/2\). The quenched Markov property at that first visit gives
\[
P_{x,\omega}(\text{visit }A_1\text{ by the top}\mid D)
\le2P_{x,\omega}(E_1\mid D),
\]
and the analogous inequality holds for the walk from \(y\). Every common
site belongs to at least one \(A_i\). A union bound proves the contact
bound; no synchronization of the two visit times is required.

\begin{lemma}[Off-diagonal comparison]\label{gauss-off-diagonal}
Fix \(a,t>0\). Uniformly over same-layer starts with
\(|y_2-x_2|\ge ar\), the total variation distance between the pair
path-prefix laws through level \(k=\lfloor tn(r)\rfloor\), under
\(P_{x,y}^{++}\) and under independent conditioned walks, satisfies
\[
\limsup {\rm dist}_{\rm TV}\le C\exp(-ca^2/t).
\tag{17}\label{eq:17}
\]
The constants depend only on the walk law.
\end{lemma}

\begin{proof}
In the top layer choose the windows
\[
I_1=\{z:|z_2-x_2-k\theta_r|<ar/3\},\qquad
I_2=\{z:|z_2-y_2-k\theta_r|<ar/3\}.
\]
They are disjoint. Lemma~\ref{gauss-shared-marginals} and the normal
tail bound show that the shared conditioned probability of either
endpoint error has limsup at most \(Ce^{-ca^2/t}\), uniformly in the
starts. Average the contact bound under the environment
law with density \(q(x)q(y)/p_2(x,y)\). This gives the same bound for
trace contact under \(P_{x,y}^{++}\).

To compare laws off contact, first use raw successful prefixes: both
walks reach level \(k\) before dropping below the common base. For any
two disjoint finite traces, their raw prefix weights agree after
averaging in shared and independent environments, because their
departure-row sets are disjoint. Passing from two finite successes to
two infinite no-drop events removes raw mass at most \(2(u_k-p)\)
under either pair law. Moreover \(p_2(x,y)-p^2\to0\) uniformly for the
present starts, since their separation tends to infinity. Thus the
conditioned prefix laws agree off contact up to a variation error that
tends to zero. Their total masses are both one, so the contact mass on
the independent side has the same upper bound up to this error. The
contact estimate proves \eqref{eq:17}.
\end{proof}

\subsection{Zero occupation on the projected diagonal}

The preceding comparison applies when the second-coordinate gap is of
order \(r\). It does not exclude a limiting pair that spends positive
time at zero gap. We now bound that occupation. The estimate counts
common true cuts first; exact-truth factorization will then transfer
it to every deterministic layer.

\begin{proposition}[Vanishing diagonal occupation]\label{gauss-diagonal}
Along every Gaussian sequence, under \(P_{0,0}^{++}\), for each fixed
finite \(T\),
\[
\lim_{\rho\downarrow0}\limsup_{r\to\infty}
\frac1{n(r)}E_{0,0}^{++}
 \sum_{0\le h\le Tn(r)}\mathbf1_{\{|Z_h|\le\rho r\}}=0.
\tag{18}\label{eq:18}
\]
\end{proposition}

\begin{proof}
Choose \(1<\alpha<\gamma<2\). The potential will be a flattened and
capped version of \(u^\alpha\), where \(u\) is the absolute gap at an
observed common cut. Its expected increase will pay for the number of
near-diagonal cuts before the next observation. We need one estimate
at nearly Gaussian radii and another below the largest remaining
exceptional radius.

\paragraph{Drift at nearly Gaussian radii.}
Fix a sufficiently small \(t_0>0\). There exist \(l\in(0,1)\),
\(\epsilon>0\), and \(u_0<\infty\) such that every radius satisfying
\[
u\ge u_0,\qquad n(u)\Pr(|S|>lu)\le\epsilon
\]
has, uniformly over shared starts with \(|Z_0|=u\),
\[
\begin{split}
E^{++}\big[(|Z_{\sigma(k(u))}|/u)^\alpha\wedge2^\alpha\big]
 &\ge1+c_*t_0,\\
P^{++}\big(\min_{0\le h\le k(u)}|Z_h|\le u/2\big)
 &\le C_*e^{-c_*/t_0},
\end{split}
\qquad k(u)=\lfloor t_0n(u)\rfloor\ge1.
\tag{19}\label{eq:19}
\]
Call these radii good. The constants \(c_*,C_*\) can be fixed
independently of small \(t_0\); \(l,\epsilon,u_0\) may depend on it.

Here is the compactness argument giving this uniform assertion. Along
any Gaussian sequence of radii \(u\), the independent gap at height
\(k(u)\), divided by \(u\), converges to \(\pm1+G\), where
\(G\sim N(0,2\sigma_*^2t_0)\). For small \(t_0\),
\[
E[|1+G|^\alpha\wedge2^\alpha]\ge1+ct_0.
\]
Indeed \(|z|^\alpha\) has strictly positive second derivative near
one, while the Gaussian probability of leaving that neighborhood is
exponentially small in \(1/t_0\). The off-diagonal comparison changes
the expectation by at most \(Ce^{-c'/t_0}\) in the limit. Replacing
\(k(u)\) by \(\sigma(k(u))\) costs \(o(1)\): by \eqref{eq:6} the
unscaled gap difference is uniformly tight, and the tested function is
bounded and uniformly continuous. No moment of the overshoot is used.

For the second estimate, shrinking the gap by \(u/2\) requires at least
one centered marginal to move by \(u/4\). The shared marginal limit
and the Brownian maximal bound
\[
\Pr\left(\sup_{s\le t_0}|B(s)|\ge x\right)
\le2\exp\left(-\frac{x^2}{2\sigma_*^2t_0}\right)
\]
give the required limit bound. This maximal bound follows by applying
the nonnegative submartingale inequality to the two exponentials of
Brownian motion on finite grids and then using continuity.

Choose strict slack in both limiting inequalities. If no triple
\((l,\epsilon,u_0)\) gave \eqref{eq:19}, choose violating radii while
\(l,\epsilon\downarrow0\) and \(u_0\to\infty\). These radii would
form a Gaussian sequence, contradicting the preceding limits. Finally,
decrease \(\epsilon\) until \(4/(1+4\epsilon)\ge2^\gamma\), and
increase \(u_0\) until \(k(u)\ge1\).

\paragraph{A cutoff above all exceptional radii.}
Fix \(A>1\). Along the given Gaussian sequence \(r\), the nongood
radii in \([u_0,Ar]\) are all \(o(r)\). For every fixed \(\xi>0\),
this follows uniformly on \([\xi r,Ar]\) from the Gaussian tail tests,
monotonicity of \(n\), and \(n(Ar)\le A^2n(r)\).

On any subsequence pass further as follows. If the nongood radii stay
bounded, choose a fixed cutoff \(b_\circ\) above them and above
\(u_0\). Otherwise choose a nongood radius \(D_*\to\infty\) within
a factor two of their supremum, and set \(b_\circ=4D_*\). In either
case \(b_\circ=o(r)\), every radius between \(b_\circ\) and \(Ar\)
is good, and
\[
\frac{n(u)}{n(r)}\le C_A(u/r)^\gamma
\qquad(b_\circ\le u\le Ar).
\tag{20}\label{eq:20}
\]
For this last bound, at a good radius
\[
m(2u)\le m(u)+4u^2\Pr(|S|>u)
\le(1+4\epsilon)m(u),
\]
so \(n(2u)\ge2^\gamma n(u)\). Iterate while the doubling stays
below \(r\). For \(r\le u\le Ar\), monotonicity of \(m\) gives
the remaining estimate. The constant \(C_A\) need not depend on
\(t_0\).

\paragraph{Escape from the cutoff region.}
Uniformly over \(|Z_0|\le b_\circ\), there is a positive probability
of reaching an exact common true level \(h\le K\) with gap at least
\(2b_\circ\), where \(K\le Cn(b_\circ)\). The constants are uniform
along the chosen further sequence; when \(b_\circ\) is fixed they
may depend on that fixed value.

For a fixed cutoff, bounded uniformly elliptic scripts to level one
can enlarge the gap beyond \(2b_\circ\); fresh joint truth completes
the claim. Consider therefore \(b_\circ=4D_*\) with \(D_*\to\infty\).
Since \(D_*\) is nongood, the second line of \eqref{eq:13}, at
\(H=\lfloor t_1n(D_*)\rfloor\), supplies one deterministic sign of
single-path deviation from \(b_H\), of size at least \(c_0D_*\) and
\(P^+\)-probability at least \(c_1t_1\). Here \(t_1>0\) is fixed
and sufficiently small, while \(c_0,c_1\) do not depend on this small
choice. The raw quenched mass of that deviation therefore exceeds
\(pc_1t_1/2\) on an environment event of probability at least
\(pc_1t_1/2\).

Assign this deviation to the rightmost walk for a positive sign, or
to the leftmost walk for a negative sign. For the other walk,
\eqref{eq:12} supplies fixed positive raw mass within
\(c_0D_*/3\) of \(b_H\), except on an environment event of probability
\(O_{c_0}(t_1^2)\). Subtracting that error from the preceding
\(\Omega(t_1)\) probability gives a uniformly positive probability
that both masses are available in the same environment. Multiplying
the quenched masses then gives a uniformly positive raw probability
of two successes whose absolute gap has increased by at least
\(c_0D_*/2\). Their prefixes use only rows below the top. Fresh joint
truth there has probability at least \(\underline p_2\), so the top
is an exact common cut and both initial no-drop events hold.

Repeat a fixed number of these blocks, choosing the outward-moving
walk at each new common cut. The shared continuation rule applies at
every repetition. Enough repetitions enlarge any initial gap in
\([0,4D_*]\) beyond \(8D_*\), with uniformly positive probability,
and their total height is a fixed multiple of \(H\). This proves the
claim with \(K\le Cn(b_\circ)\).

\paragraph{Observation times and potential increase.}
Starting at the common cut zero, observe only the following subset of
common cuts. Stop when an observed gap is at least \(Ar\). From an
observed gap \(u<Ar\), choose the next observation by the rule
\[
\begin{cases}
\text{first common cut at offset at least }k(u),&u>b_\circ,\\
\text{first subsequent common cut with gap at least }2b_\circ
\text{ or offset at least }K,&u\le b_\circ.
\end{cases}
\]
These are stopping indices of the common-word chain. At each observation
the conditional continuation is the fresh shared conditioned pair law
from the observed sites.

Use the bounded potential
\[
\Phi(u)=[(u\vee b_\circ)\wedge2Ar]^\alpha.
\]
For \(b_\circ<u<Ar\), \eqref{eq:19} implies an expected increase at
least \(c_*t_0u^\alpha\). Flattening below \(b_\circ\) only increases
the terminal value, and the upper cap is at least \(2u\). For
\(u\le b_\circ\), the potential cannot decrease; the escape event
just proved makes it increase by at least
\((2^\alpha-1)b_\circ^\alpha\) with a uniformly positive probability.
Telescope up to observed exit, first with a deterministic truncation
of the number of observations. Since \(\Phi\le(2Ar)^\alpha\),
monotone convergence gives
\[
E^{++}\sum_{\substack{\text{pre-exit observations}\\u>b_\circ}}
 t_0u^\alpha\le C(Ar)^\alpha,
\qquad
E^{++}\sum_{\substack{\text{pre-exit observations}\\u\le b_\circ}}
 b_\circ^\alpha\le C_\circ(Ar)^\alpha.
\]
The first constant is independent of small \(t_0\).

\paragraph{Counting all common cuts.}
Allocate each common cut to the last observation at or before it,
excluding the next observation. An upper-regime observation at gap
\(u\) receives at most \(k(u)\) cuts: each allocated offset is an
integer strictly below \(k(u)\). A lower-regime observation receives
at most \(K\) cuts. These bounds do not require any bound on the
height overshoot of the next observation.

Let \(N_{\rho}^{\rm pre}\) count allocated common cuts with gap at
most \(\rho r\), before observed exit. Split observations into three
sets. For fixed \(\rho>0\), take \(r\) large enough that
\(b_\circ\le2\rho r\). Equation~\eqref{eq:20} and the drift bounds give
\[
\begin{array}{c|c}
\text{gap at the observation}&
\text{contribution to }E^{++}N_{\rho}^{\rm pre}/n(r)\\[2pt]\hline
u\le b_\circ&C_{A,\circ}(b_\circ/r)^{\gamma-\alpha}\\
b_\circ<u\le2\rho r&C_A\rho^{\gamma-\alpha}\\
2\rho r<u<Ar&C_Ae^{-c_*/t_0}.
\end{array}
\]
For the first row use \(K\le Cn(b_\circ)\). For the second use
\(k(u)\le t_0n(u)\) and
\(u^{\gamma-\alpha}\le(2\rho r)^{\gamma-\alpha}\). For the last
row, an allocated near-diagonal cut requires shrinking the gap from
\(u\) to at most \(u/2\) before offset \(k(u)\). Its conditional
probability is bounded by \eqref{eq:19}; the number of allocated
cuts is at most \(k(u)\). In both upper-regime rows, the factor \(t_0\)
from \(k(u)\) cancels the one in the drift budget. Their constants
can therefore be independent of small \(t_0\).

Under the same-site initial law, joint marginal tightness makes the
probability of an observed gap \(\ge Ar\) before height \(Tn(r)\)
tend to zero as \(A\to\infty\), uniformly in the limiting upper
bound in \(r\). All common cuts after that exit and below the target
height contribute at most \(Tn(r)+1\) times its indicator.

\paragraph{From common cuts to deterministic layers.}
Let \(\mathcal C_h\) denote common truth at height \(h\), and let
\(A_h\) be the raw probability that both paths reach \(h\) before
the initial drop with \(|Z_h|\le\rho r\). Imposing the initial
shared conditioning gives a numerator at most \(A_h\). Instead,
impose both fresh truths at height \(h\). Their conditional probability
is at least \(\underline p_2\), and the resulting event implies the
initial conditioning and \(\mathcal C_h\). Dividing by the same
initial normalizer gives
\[
P_{0,0}^{++}(|Z_h|\le\rho r)
\le\underline p_2^{-1}
P_{0,0}^{++}(|Z_h|\le\rho r,\mathcal C_h).
\]
Sum over \(h\le Tn(r)\) and use the common-cut bounds above. On the
chosen further sequence the lower-regime term vanishes because
\(b_\circ=o(r)\). The remaining pre-exit bound is
\(C_A\rho^{\gamma-\alpha}+C_Ae^{-c_*/t_0}\). Every subsequence
admits this further-sequence construction, so the bound controls the
original limsup. Send \(\rho\downarrow0\), then \(t_0\downarrow0\),
and finally \(A\to\infty\). This proves \eqref{eq:18}.
\end{proof}

\subsection{Shared joint limit and quenched consequence}

\begin{proposition}[Independent joint limit for shared paths]\label{gauss-joint}
Under \(P_{0,0}^{++}\), the pair of processes in \eqref{eq:16}
converges functionally to two independent Brownian motions, each with
variance parameter \(\sigma_*^2\).
\end{proposition}

\begin{proof}
Take any subsequential continuous-path joint limit \(B_1,B_2\) of \eqref{eq:16} under \(P_{0,0}^{++}\), on fixed horizons (enlarge a horizon if needed at endpoints), and let \(\mathcal G_s\) be the joint past. By \eqref{eq:18}
\[
\int_0^T\mathbf1_{\{B_1(s)=B_2(s)\}}\,ds=0\quad\text{almost surely}.
\tag{21}\label{eq:21}
\]
One can integrate continuous cutoffs of the gap supported in shrinking neighborhoods of zero; their grid sums give the same limiting expectations by joint tightness with continuous limits.

\emph{Conditional marginal increments.}
At a deterministic time \(s\), put \(H_r=\lfloor sn(r)\rfloor\)
and restart the prelimit pair at \(\sigma(H_r)\). This is a stopping
index of the chain of common words, not a stopping time for an
individual walk. All first-hit coordinates of both paths at levels
at most \(H_r\) are determined by the words traversed through this
boundary. Its conditional suffix law is the fresh shared conditioned
law from the terminal pair of sites.

By \eqref{eq:6}, \(\sigma(H_r)-H_r\) is uniformly tight.
Consequently its normalized time shift vanishes. Joint tightness with
continuous limits, on a slightly enlarged horizon, makes the values
at \(H_r+\lfloor tn(r)\rfloor\) and at
\(\sigma(H_r)+\lfloor tn(r)\rfloor\) differ by \(o_{\Pr}(r)\)
after subtracting the same linear centering. At the initial endpoints
this also follows directly from \eqref{eq:6}; the centering correction
is a tight integer times \(\theta_r/r\), which vanishes by
Lemma~\ref{gauss-medians}.

Test against bounded continuous functions of finitely many joint past
coordinates, represented at floor grid times in the prelimit. Condition
at the common boundary and test the fresh normalized increment of
length \(\lfloor tn(r)\rfloor\). Lemma~\ref{gauss-shared-marginals}
makes its conditional bounded continuous test expectation converge to
the normal value uniformly in the random terminal pair of sites.
The time replacements just justified do not change the limit. Passing
to the joint limit, then using a monotone-class argument for the joint
past, proves that for each \(t>0\),
\[
\mathcal L(B_i(s+t)-B_i(s)\mid\mathcal G_s)
=N(0,\sigma_*^2t),\qquad i=1,2.
\]

\emph{Off-diagonal cross increments.}
Moreover if \(g\) is a continuous \([0,1]\)-valued cutoff vanishing on \([-2a,2a]\), \(a>0\), and \(\delta_i=B_i(s+t)-B_i(s)\), then for bounded \(\mathcal G_s\)-measurable \(F\),
\[
\big|E[F g(B_2(s)-B_1(s))\delta_1\delta_2]\big|
\le \|F\|_\infty\, t\,e_a(t),\qquad e_a(t)\longrightarrow0\quad(t\downarrow0).
\tag{22}\label{eq:22}
\]
The bound can be uniform in \(s\). To see this, again restart as above; on the supported past gaps the restart state has absolute gap \(\ge a r\) except with vanishing error. Clip each future normalized increment symmetrically to the interval \([-w\sqrt t,w\sqrt t]\), for the moment fixed \(w,t\). By \eqref{eq:17} at the fresh state their product expectation differs from the independent-law one by at most \(Cw^2 t\exp(-c a^2/t)\) in limsup; the independent-law value tends to zero by the single Brownian limits. Test first against continuous bounded past functions and pass to the joint limit. Removing the clipping there costs at most \(t\,o_{w\to\infty}(1)\|F\|_\infty\), by the conditional marginal normals and Cauchy--Schwarz. Continuous past tests extend to bounded measurable ones, for instance by bounded approximation using \(E(|\delta_1\delta_2|\mid\mathcal G_s)\le \sigma_*^2 t\). Choosing \(w\to\infty\) slowly gives \eqref{eq:22}.

\emph{Identification of the joint law.}
These facts force \(B=(B_1,B_2)\) to consist of independent Brownian motions. For a direct verification fix \(s_0\), bounded \(F\) of the joint past there and a real vector \(\xi\in\mathbb R^2\). Telescope the exponential of \(i\xi\cdot(B(s)-B(s_0))\) on an equal mesh of \([s_0,v]\), multiply by \(F\) and take expectation. Taylor remainders through order two total \(o(1)\) by the marginal Gaussian third absolute moment bounds. First-order terms vanish by the past-conditional laws, and diagonal second-order terms give the Riemann sum using their variances. Cross terms total \(o(1)\): away from gap zero by \eqref{eq:22}, and the complementary terms are bounded by a constant times the mesh sum of step lengths times expected continuous cutoff near gap zero (by the past-conditional absolute product bound), which is negligible by continuity and \eqref{eq:21} as the cutoff shrinks. Thus \(f(v)=E[F\exp(i\xi\cdot(B(v)-B(s_0)))]\) satisfies
\[
f(v)=E F-\tfrac12\sigma_*^2|\xi|^2\int_{s_0}^v f(s)\,ds .
\]
This gives the Gaussian increment law independent of the joint past, proving the assertion.
\end{proof}

The passage from two copies to quenched concentration is a second-moment
method developed in RWRE by Bolthausen and
Sznitman~\cite[Lemma~4.1]{BolthausenSznitman2002} and Berger and
Zeitouni~\cite[Section~2]{BergerZeitouni2008}. The conclusion needed here
is convergence in environment probability along Gaussian sequences; the
following calculation supplies it under the present hypotheses.

\begin{corollary}[Quenched Gaussian mass]\label{gauss-quenched}
Along Gaussian sequences, the quenched path under \(D\)-conditioning
from a fixed start has the Brownian limit of \eqref{eq:16} in
\(P\)-probability, in the sense of bounded continuous path tests.
In particular, for every \(\rho>0\) and every \(b_*(r)\to\infty\),
with \(H=\lfloor n(r)\rfloor\),
\[
P\left(\frac{K_H(0,\{ W_H\le b_*(r)r,\ |Y_H-b_H|\le\rho r\})}{a_H}
\ge c_\rho\right)\longrightarrow1
\tag{23}\label{eq:23}
\]
for some \(c_\rho>0\) depending only on \(\rho\) and the walk law.
The same assertion holds at translated starts.
\end{corollary}

\begin{proof}
Let \(F\) be a bounded real continuous test on path space, and let
\(M_r(\omega)\) be its expectation under the quenched conditioned
path. Give the environment the probability law
\[
\widehat P(d\omega)=\frac{q(0)^2}{E[q(0)^2]}P(d\omega).
\]
Under this weighting, two conditionally independent quenched
conditioned paths have exactly the annealed law \(P_{0,0}^{++}\).
Proposition~\ref{gauss-joint} and its marginal limits therefore give,
with \(b=E[F(B)]\) for the limiting Brownian path,
\[
E_{\widehat P}M_r\longrightarrow b,
\qquad E_{\widehat P}M_r^2\longrightarrow b^2.
\]
Thus \(M_r\to b\) in \(\widehat P\)-probability. Since
\(q(0)>0\) almost surely, \(P\) is absolutely continuous with respect to
\(\widehat P\); convergence in probability also holds under \(P\).

For \eqref{eq:23}, use Lemma~\ref{gauss-medians} to replace the
linear centering by the medians. Choose a nonnegative continuous
path test below the desired event, with a fixed bounded supremum and
with the endpoint, throughout a small neighborhood of time one,
strictly inside the window of radius \(\rho\). Such a test can be
chosen with positive Brownian expectation: the time-one centered normal
law assigns positive mass to a smaller window, and continuity permits
the nearby-time and sufficiently large fixed supremum restrictions.
Eventually this fixed supremum bound is smaller than \(b_*(r)\).
The quenched no-drop-conditioned probability is consequently bounded
below by a fixed positive constant with \(P\)-probability tending to one.
Finally, on successful prefixes the total variation cost of replacing
no-drop conditioning by conditioning on \(T_H<T_{-1}\) is at most
\(1-q(0)/a_H\), which tends to zero almost surely. This proves the
stated normalization by \(a_H\). Translation invariance gives the
same conclusion at any translated start.
\end{proof}
